\subsection{Experimental Setup and Evaluation Protocol}

This section outlines the experimental methodology that will be used to evaluate the proposed Audio Deepfake Detection (ADD) framework under realistic RTC communication environments. The evaluation protocol is designed to assess both detection effectiveness and deployment feasibility under various transmission conditions commonly encountered in VoIP and RTC systems.

The experiments will focus on three primary aspects: robustness against RTC-induced distortions, comparative analysis of multiple detection architectures, and assessment of computational feasibility for real-time deployment.

\subsubsection{Dataset Configuration}

The experimental evaluation relies on two dataset components: the ASVspoof 5 benchmark and a custom RTC-aware augmentation pipeline applied on top of it.

\textbf{ASVspoof 5 (2024/2026):} Focuses on ``Generic Telephony or VoIP'' scenarios. It incorporates crowdsourced data recorded in diverse acoustic conditions and processes it with both conventional and contemporary neural codecs (e.g., EnCodec) to simulate modern VoIP quality.

\textbf{RTC-Aware Augmentation Pipeline:} Applied on top of ASVspoof 5 audio, this custom pipeline introduces codec round-trip transformation, bursty packet-loss simulation, and receiver-side post-processing to expose training and evaluation data to RTC-induced transmission conditions beyond those captured in ASVspoof 5 alone.

The dataset is partitioned into three speaker-disjoint subsets:

\textbf{Training Set:} consist balanced real/fake utterances augmented with simulated codecs and packet loss

\textbf{Development Set:} used for hyperparameter tuning and setting detection thresholds

\textbf{Evaluation Set:} will include unseen spoofing algorithms to rigorously test the system’s cross-domain generalization and real-time performance

\subsubsection{Synthetic Speech Generation Methods}

The evaluation datasets used in this study contain synthetic speech generated using multiple state-of-the-art speech synthesis techniques. These techniques represent the primary attack categories encountered in modern audio deepfake systems.

The synthetic speech samples include:

\begin{itemize}
\item \textbf{Text-to-Speech (TTS):} Systems that generate speech directly from text input using neural speech synthesis architectures.

\item \textbf{Voice Conversion (VC):} Systems that transform the voice characteristics of a source speaker into those of a target speaker while preserving linguistic content.

\item \textbf{Neural Voice Cloning:} Systems that synthesize speech resembling a target speaker using a limited amount of reference audio.

\item \textbf{Modern Generative Paradigms:} Recent advancements utilize Diffusion Models (e.g., Grad-TTS) and Neural Vocoders (e.g., HiFi-GAN, BigVGAN) to achieve high-fidelity, efficient waveform synthesis that is increasingly difficult for human listeners to distinguish from real speech
\end{itemize}

The selected datasets contain speech generated using a variety of synthesis architectures, including autoregressive, non-autoregressive, diffusion-based, and neural codec-based speech generation models.

\subsubsection{RTC Distortion Robustness Evaluation}

One of the primary objectives of this research is to investigate the impact of RTC communication channels on deepfake detection performance. Therefore, the proposed framework will be evaluated under multiple transmission conditions representing common VoIP environments.

The evaluation will consider variations in:

\begin{itemize}
\item Codec compression using Opus, G.711, SILK, and AMR-WB.
\item Packet loss rates of 0\%, 5\%, 10\%, 15\%, and 20\%.
\item Network delay and jitter conditions representing stable and unstable communication channels.
\item Bandwidth-constrained scenarios with adaptive bitrate operation.
\item Combined impairment scenarios involving simultaneous codec compression, packet loss, and network variability.
\end{itemize}

By systematically varying these parameters, the study aims to determine the robustness of the proposed detection framework against realistic communication impairments.


\subsubsection{Comparative Model Evaluation}

The proposed study will investigate multiple deepfake detection architectures to identify the most suitable model for deployment in RTC environments.

Candidate models include:

\begin{itemize}
\item AASIST
\item XGBoost-based classifiers
\end{itemize}

In addition, different feature extraction approaches will be explored using self-supervised speech representations such as WavLM and Wav2Vec~2.0.

All models will be evaluated using the performance metrics defined in Section 4.7. Comparative analysis will be performed to examine differences in detection accuracy, robustness to transmission distortions, and suitability for real-time deployment.

\subsubsection{Computational Feasibility Assessment}

Since the proposed framework is intended for integration within real-time communication systems, computational efficiency is an important consideration.

The experimental study will therefore include an assessment of runtime characteristics, including:

\begin{itemize}
\item Inference latency
\item Throughput under concurrent audio streams
\item Memory consumption
\item CPU and GPU utilization
\item Model size and deployment complexity
\end{itemize}

This analysis will help determine whether the proposed framework can satisfy the latency and scalability requirements of practical RTC applications.

\subsubsection{Evaluation Objectives}

The experimental protocol is designed to answer the following research questions:
\vspace{-0.5cm}
\begin{enumerate}[label=\textbf{\arabic*.}]
\item How do RTC transmission distortions affect audio deepfake detection performance?
\item Which combination of feature extraction and classification architecture provides the highest robustness under RTC conditions?
\item What trade-offs exist between detection accuracy and computational efficiency?
\item Can the proposed framework satisfy the requirements of real-time deployment within RTC systems?

\end{enumerate}